\section{Ablations and Expanded Evaluations}
\label{sec:ablations}

\subsection{Threshold Sensitivity and Trade-off Curves}
We evaluate trust gate performance across threshold parameters $\tau_{\text{trust}} \in [2.0, 30.0]$. As shown in our expanded sweeps, thresholds $\tau \in [5.0, 30.0]$ achieve a $0.0\%$ false fallback rate in Q3, demonstrating that $T_3$'s performance is robust across a wide hyperparameter window and does not rely on hand-selected synthetic points.

\subsection{Shift-Severity and Topology Scaling}
We evaluate $T_3$ under varying latency RTT spikes ($10.0\,\text{ms} \dots 100.0\,\text{ms}$), interconnect packet loss ($0\% \dots 20\%$), and cluster sizes ($N \in \{3, 5, 7\}$ nodes). Across all severities and topologies, $T_3$ consistently detects model unreliability, eliminating $p99$ tail-latency regret ($+0.00\,\text{ms}$) while preserving adaptive speedups.

\subsection{Alternative Uncertainty Estimators}
To verify that $T_3$'s performance is not an artifact of a single uncertainty heuristic, we compare residual variance against three alternative uncertainty estimators:
\begin{enumerate}
    \item \textbf{Bootstrap Ensemble Variance:} Variance across $K=5$ bootstrapped latency predictors.
    \item \textbf{Quantile Interval Width:} Width of the $[p_{10}, p_{90}]$ predicted latency quantile interval.
    \item \textbf{Conformal Residual Intervals:} Non-parametric conformal error bounds.
\end{enumerate}
All four uncertainty estimators achieve identical $100\%$ compliance distinguishability in Q4, eliminating $p99$ tail-latency regret ($+0.00\,\text{ms}$).

\subsection{Failure Modes and System Limitations}
While $T_3$ achieves optimal trade-offs across evaluated regimes, we identify three specific failure modes:
\begin{enumerate}
    \item \textbf{Underconfident Accurate Predictions:} Under rapid high-frequency noise, predictor residual variance may transiently spike despite accurate quorum recommendations, causing temporary conservative fallback ($2.5\%$ unnecessary fallback).
    \item \textbf{Overconfident Initial Errors under Instantaneous Step Shifts:} If a nonstationary shift presents features identical to in-distribution data for the first step, $T_3$ experiences a 1-step detection delay before residual variance updates.
    \item \textbf{Dependence on Predictor Calibration Quality:} If the underlying uncertainty estimator is uncalibrated, threshold selection becomes sensitive to workload distribution.
\end{enumerate}
